\subsection*{4. Structure of deficient triples}
\addcontentsline{toc}{subsection}{4. Structure of deficient triples}
\paragraph{Theorem G1.} Let $T\in H$ be deficient in an avoiding $3$-graph. Then:
\begin{enumerate}
\item[(i)] $T$ has exactly one triangle option $q=xy$, with $x\in T$. Write $\CN(q)=\{uv,vw,uw\}$ and $T=x+uv$.
\item[(ii)] Exactly one of $xu,xv$ is light.
\item[(iii)] If $xu$ is light, then $L(xu)=\{v,w\}$ and $s_{vw}=2$. If $xv$ is light, then $L(xv)=\{u,w\}$ and $s_{uw}=2$.
\end{enumerate}
Moreover $T$ has no isolated slot, and its precise slot structure is given by Corollary 6g below. The inputs are F1, the slot and profile lemmas of subdivision 2, the profile theorem there, and Lemma 3g. No redistribution conclusion is used.

\subsubsection*{4.1 Translation into profiles}
\paragraph{Lemma 1g.} Fix $T=abc$ in an avoiding $H$.
\begin{enumerate}
\item[(a)] Triangle options are exactly the options of type-2 slots. Distinct slots give distinct options. Thus $T$ has $t$ triangle options.
\item[(b)] If $(z,A)$ is type~2, with partner slots $(w,B),(w,C)$, then $\CN(a,z)=\{bc,bw,cw\}$. The pair $ab$ is light exactly when $|C|=1$, in which case $C=\{w\}$ and $L(ab)=\{c,w\}$. The corresponding assertion for $ac,B$ also holds.
\item[(c)] If $t=1$, then $(w,B)$ is type~1 in $R_c$ and $(w,C)$ is type~1 in $R_b$. Hence $bc,cb\geq1$, where these symbols denote profile counts.
\item[(d)] If $t=1$ and $C=\{w\}$, then $s_{cw}=2+ab$ and the option through the pair $ab$ is $cw$. Symmetrically, if $B=\{w\}$, then $s_{bw}=2+ac$.
\end{enumerate}
All assertions rotate cyclically with $a,b,c$.
\begin{proof}
For (a), the members of $\CN(a,z)$ are $bc$, the pairs $bu$ coming from $R_b$-neighbours, and the pairs $cu$ coming from $R_c$-neighbours. Lemma 2 gives the triangle $\{bc,bw,cw\}$ for a type-2 slot. An isolated slot gives just one pair. A type-1 slot gives pairs with a common vertex, hence cannot give a triangle. Options of distinct slots are distinct because an option meets $T$ only at its opposite vertex.
For (b), $k_{ab}=1+|C|$ and $w\in C$; similarly for $ac$. For (c), each partner slot has its edge to $(z,A)$, and cannot be type~2 because that would give a second type-2 slot.
For (d), the slot $(w,C)$ has as $R_b$-neighbours the type-2 slot $(z,A)$ and every slot counted by $ab$: a type-1 slot in $R_b$ needs a neighbour, and $C$ consists only of $w$. It has no $R_a$-neighbour, since it is not type~2. Its degree is $1+ab$, giving $s_{cw}=2+ab$ by Lemma 1. The mirror assertion is identical after exchanging $B,C$.
\end{proof}
A feasible profile with $t=1$ and its type-2 slot in $X_i$ is \emph{G1-compliant} when exactly one of $n_{i+1}=1,n_{i-1}=1$ holds and
\[
 n_{i-1}=1\Longrightarrow g_i=0,\qquad n_{i+1}=1\Longrightarrow f_i=0.
\]
Its \emph{partner constraints} are $g_{i+1}\geq1$ and $f_{i-1}\geq1$. These follow already from (F) when $t=1$. All these notions are invariant under cyclic shift.
\paragraph{Corollary 1g.} If $T$ has exactly one type-2 slot and no isolated slot, its profile satisfies the partner constraints, and conclusions (i)--(iii) of G1 are equivalent to G1-compliance.
\begin{proof}
Rotate to $i=0$. Lemma 1g(c) gives the partner constraints. Without isolated slots, $|B|=n_B$ and $|C|=n_C$. By Lemma 1g(b), exactly one light pair means exactly one of $n_B,n_C$ equals one. By Lemma 1g(d), the required common-link size two then says $ab=g_0=0$ if $C=\{w\}$, or $ac=f_0=0$ if $B=\{w\}$. The link identity in G1(iii) is part of Lemma 1g(b).
\end{proof}
\subsubsection*{4.2 Isolated slots}
\paragraph{Lemma 2g.} A triple with exactly one type-2 slot and at least one isolated slot satisfies $\sigma(T)+\tau(T)\geq9/4$.
\begin{proof}
Rotate its type-2 slot to $(z,A)$. Each group has at least one non-isolated slot, by Lemma 1g(c). For $n\geq1$, define
\[
 D(i)=\frac1{1+n+i}+\frac i2-\frac1{1+n},\qquad i\geq0.
\]
Then $D(0)=0$ and
\[
 D(i+1)-D(i)=\frac12-\frac1{(1+n+i)(2+n+i)}\geq\frac13.
\]
One isolated slot therefore supplies a correction at least $1/3$ in Lemma $3^+$. Theorem 2 gives $\sigma+\tau\geq23/12+1/3=9/4$.
\end{proof}
\subsubsection*{4.3 The threshold-two profile theorem}
\paragraph{Theorem 4g.} For every feasible profile: (a) $t\geq2$ implies $\Phi_{\rm prof}\geq2$; (b) if $t=1$, the partner constraints hold, and the profile is not G1-compliant, then $\Phi_{\rm prof}\geq2$.
Equality in (a) occurs at the complete five-vertex profile $t_A=t_B=t_C=2$, with all type-1 counts zero. Only the bound two is asserted in (b).
\begin{proof}
The bound $\Phi_{\rm prof}\geq t/6$ disposes of $t\geq12$, so take $1\leq t\leq11$. A big relation again has both type-1 counts at least 23.
For at least two big relations, the estimate in Theorem 2 gives
\[
 \Phi_{\rm prof}\geq2+t\left(\frac16-\frac1{576}\right)-\frac1{12}
 \geq2+\frac{95}{576}-\frac{48}{576}>2.
\]
For a non-big relation with one count at least 93, (E2) gives $\Phi_{\rm prof}\geq8649/4300>2$. Thus every count of a non-big relation can be restricted to $0,\ldots,92$.
If no relation is big, certificate C1g proves (a) for the 360 triples of type-2 counts satisfying $2\leq t\leq11$. Exactly 41 feasible profiles attain two. For $t=1$, let $i$ be the position of the type-2 slot. Under the partner constraints, the noncompliant profiles are precisely the union of
\begin{align*}
 V1:&\quad n_{i+1}\geq2,\quad n_{i-1}\geq2;\\
 V2:&\quad n_{i+1}=n_{i-1}=1;\\
 V3:&\quad n_{i-1}=1,\quad g_i\geq1;\\
 V4:&\quad n_{i+1}=1,\quad f_i\geq1.
\end{align*}
Indeed both neighbouring group sizes are at least one, so these are exactly the negation of G1-compliance. Each class is expressible as masks on the three group matrices of the profile cycle. The C1g computation minimises over those masks, for all three positions, and gives the bound two in every class.
If exactly one relation is big, rotate it to $R_c$. The notation of Theorem 2 gives
\[
 \Phi_{\rm prof}\geq1+W,\qquad
 W=G+\kappa^*(ab+t_A,t_A)+\kappa^*(ba+t_B,t_B),
\]
where $ab,ba,ca,cb\in\{0,\ldots,92\}$ and $n_A,n_B\geq23$. For $t=1$ the compliance alternatives are:
\begin{itemize}
\item $(t_A,t_B,t_C)=(1,0,0)$: compliance means $n_C=ca+cb=1$ and $ab=0$; the partner constraint is $cb\geq1$.
\item $(t_A,t_B,t_C)=(0,1,0)$: compliance means $n_C=1$ and $ba=0$; the partner constraint is $ca\geq1$.
\item $(t_A,t_B,t_C)=(0,0,1)$: compliance is impossible since both neighbouring groups have more than one slot.
\end{itemize}
Certificate C2g gives $W\geq1$ for $2\leq t\leq11$ and for every noncompliant $t=1$ rest profile, with or without the partner constraints. There are eight target-equality profiles in the former range. Hence $\Phi_{\rm prof}\geq2$. The cases are exhaustive, including the earlier reduction of every large one-sided count.
\end{proof}
\subsubsection*{4.4 G1 and the precise slot structure}
\begin{proof}[Proof of Theorem G1]
For deficient $T$, Lemma $3^+$ gives
\[
 \Phi_{\rm prof}\leq\Phi_{\rm prof}^+\leq\sigma(T)+\tau(T)<2.
\]
Lemma 3g forces $t\geq1$ and Theorem 4g(a) forces $t\leq1$. Thus $t=1$. Lemma 2g excludes isolated slots. Corollary 1g supplies the partner constraints, Theorem 4g(b) forces compliance, and Corollary 1g translates compliance into G1(i)--(iii).
\end{proof}
The strict hypothesis is necessary: every triple of the complete five-vertex $3$-graph has $\sigma+\tau=2$ and six triangle options. Exactly one triangle option is equivalent to membership in exactly one bicone.
\paragraph{Corollary 6g.} Suppose $T=abc$ is deficient, its type-2 slot is $(z,A)$ with partner $w$, and $C=\{w\}$, so the pair $ab$ is light. Then
\[
 t_A=1,\quad t_B=t_C=0,\quad ab=ca=ba=0,\quad cb=1,
\]
with no isolated slots. Every slot of $A$ except $(z,A)$ and every slot of $B$ is type~1 in $R_c$. The relation $R_b$ consists of the single edge $(w,C)(z,A)$, and $R_a$ is empty. In particular $|B|=bc$ and $|A|=1+ac$.
\begin{proof}
G1(iii) and Lemma 1g(d) give $ab=0$. The only slot of $C$ is type~1 in $R_b$, so $cb=1,ca=t_C=0$. As $t=1$, also $t_B=0$. Feasibility for $R_a$ forces $ba=0$. The remaining assertions follow directly from these counts and absence of isolated slots.
\end{proof}
